\section{Miasto a wieś - analiza porównawcza}

% SLAJD 16
\begin{frame}{Wyzwania życia w mieście \cite{link2}}
    \begin{itemize}
        \item \textbf{Stres i presja czasu:}
        \begin{itemize}
            \item Miasto to nieustanny wyścig z terminami i ludźmi.
            \item Nadmiar bodźców generuje chroniczny stres.
            \item \textit{Wieś:} Spokojniejszy rytm pozwala na regenerację.
        \end{itemize}
        
        \item \textbf{Koszty życia:}
        \begin{itemize}
            \item Astronomiczne ceny mieszkań ("mała przestrzeń za fortunę").
            \item \textit{Wieś:} Ten sam budżet pozwala na dom z ogrodem.
        \end{itemize}

        \item \textbf{Wszechobecny hałas:}
        \begin{itemize}
            \item Syreny, ruch uliczny, brak ciszy nocnej (pogorszony sen).
            \item \textit{Wieś:} Cisza przerywana jedynie odgłosami natury.
        \end{itemize}
    \end{itemize}
\end{frame}

% SLAJD 17
\begin{frame}{Otoczenie społeczne i przyrodnicze}
    \begin{itemize}
        \item \textbf{Deficyt natury:} Miasto to "betonowa dżungla". Brak codziennego dostępu do zieleni negatywnie wpływa na samopoczucie.
        
        \item \textbf{Przestrzeń i prywatność:}
        \begin{itemize}
            \item W bloku słychać sąsiadów przez ścianę.
            \item Wieś oferuje własne podwórko i brak tłoku.
        \end{itemize}

        \item \textbf{Relacje międzyludzkie:}
        \begin{itemize}
            \item \textit{Miasto:} Anonimowość i samotność w tłumie.
            \item \textit{Wieś:} Silne więzi sąsiedzkie i wzajemna pomoc.
        \end{itemize}

        \item \textbf{Bezpieczeństwo:} Niższy wskaźnik przestępczości na wsi i brak anonimowości sprzyjają bezpieczeństwu (szczególnie dzieci).
    \end{itemize}
\end{frame}

% SLAJD 18
\section{Tempo życia}
\begin{frame}{Różnice w codziennym tempie życia}
    \begin{columns}
        \begin{column}{0.5\textwidth}
            \textbf{\textcolor{red}{Miasto (Dyktando zegarka)}}
            \begin{itemize}
                \item Dzień precyzyjnie zaplanowany ("wypełniony po brzegi").
                \item Dominacja pośpiechu, korków i multitaskingu.
                \item Posiłki w biegu (lunch przy biurku).
                \item Wieczorem: szukanie kolejnych bodźców zamiast wyciszenia.
            \end{itemize}
        \end{column}

        \begin{column}{0.5\textwidth}
            \textbf{\textcolor{green!60!black}{Wieś (Cykl natury)}}
            \begin{itemize}
                \item Tempo wyznaczają pory roku i dnia, a nie kalendarz.
                \item Spokojniejsze poranki bez presji bycia "wszędzie naraz".
                \item Praca przeplatana odpoczynkiem (np. w ogrodzie).
                \item Wieczory spędzane z bliskimi na wyciszeniu.
            \end{itemize}
        \end{column}
    \end{columns}
\end{frame}